
\documentclass[11pt, a4paper]{article}

% ===== ESSENTIAL PACKAGES =====
\usepackage{geometry}
\geometry{a4paper, margin=1.2in}
\usepackage{setspace}
\usepackage{array}
\usepackage{xcolor}
\onehalfspacing
\setlength{\parskip}{1em}
\usepackage{multicol}

\usepackage{catchfile}

% ===== WORD COUNT COMMAND =====
\newcommand{\wordcount}{%
    \ifnum\pdf@shellescape=1%
        \immediate\write18{texcount -1 -sum Oblig1.tex > wordcount.txt}%
        \CatchFileDef{\words}{wordcount.txt}{}%
        \words%
    \else%
        [Enable -shell-escape]%
    \fi%
}
\usepackage{mdframed}

\newmdenv[
    leftmargin=-50pt,
    rightmargin=0pt,
    innerleftmargin=3pt,
    innerrightmargin=0pt,
    innertopmargin=5pt,
    innerbottommargin=5pt,
    skipabove=5pt,
    skipbelow=5pt,
    linewidth=1pt,
    linecolor=black,
    topline=true,
    rightline=false,
    bottomline=false
]{sectionboxinner}

\newenvironment{sectionbox}[1][]{%
    \begin{sectionboxinner}
    \textbf{\large #1}
    \vspace{1pt}
    \newline
}{%
    \end{sectionboxinner}
}

\newmdenv[
    leftmargin=0pt,
    rightmargin=0pt,
    innerleftmargin=0pt,
    innerrightmargin=0pt,
    innertopmargin=3pt,
    innerbottommargin=3pt,
    skipabove=8pt,
    skipbelow=8pt,
    linewidth=1.2pt,
    linecolor=gray!60,  % Lighter gray color
    topline=false,
    rightline=false,
    bottomline=false
]{subsectionboxinner}

\newenvironment{subsectionbox}[1][]{%
    \begin{subsectionboxinner}
    \textbf{#1}
    \vspace{3pt}
    \newline
}{%
    \end{subsectionboxinner}
}

% ===== CUSTOM ENVIRONMENTS FOR PHILOSOPHICAL ARGUMENTS =====
\newenvironment{argument}{\begin{quote}\itshape\textbf{Argument: }}{\end{quote}}
\newenvironment{objection}{\begin{quote}\itshape\textbf{Innvending: }}{\end{quote}}
\newenvironment{reply}{\begin{quote}\itshape\textbf{Reply: }}{\end{quote}}

% ===== DOCUMENT BEGIN =====
\title{Kap. 16 - Er moralen objektiv?}
\author{Oliver Ekeberg}
\date{\today}

\begin{document}
\maketitle

\tableofcontents

\begin{flushleft}
    Vi skal her studere om moralen er avhengig av oss, eller om den eksisterer uavhengig. Er altså moralen objektiv? Dette spørsmålet dreier seg om \textbf{\textit{metaetikk}}. Huxley argumenterer for at det vi vet om evolusjon, viser at moralen \textit{ikke} er objektiv.\newline

    Først noen viktige synspunkt i metaetikk
    \begin{enumerate}
        \item Moralteori blir beskyldt for å være tannløse. Hva vi gjør avhenger stort av hvordan vi føler oss, viser psykologisk forskning. Denne typen kritikk er rettet mot den praktiske verdien av moralsk teoretisering.
        \item Moralteori kritiseres fordi utviklere av moralteorier, var ikke moralsk prisverdige folk. Aristoteles trodde noen mennesker var født til å være slaver, og at kvinner hadde dårligere fornuftsevner enn menn. Samtidig som ideene om fellesmenneskelig menneskeverd spredte seg blant europeiske intellektuelle, begynte europeere å kolonisere store deler av verden. Dette kan bli sett på som en generalisering da. Det var ikke datidens professorer som koloniserte Asia, Afrika og Amerika, men militærledere og ekspedisjonsfolk. Vi kan bare analysere om de tok en positiv stilling til koloniseringen som skjedde. Hva er poenget med moralteori, om de som utviklet dem kunne delta og bidra til forkastelige praksiser som slaveri.
    \end{enumerate}
\end{flushleft}

\begin{sectionbox}[Synsing eller fakta?]
\begin{flushleft}
    Finnes det objektive svar på spørsmål om rett og galt? Er moral et spørsmål om fakta eller personlige meninger? Vi trenger noen begrepsmessige distinksjoner. Vi kan skille mellom ulike deler av metaetikk.\newline

    Første del av metaetikken omhandler moralens natur. Dette er spørsmål om moralske verdier finnes, og hva slags tyoe ting slike verdier eventuelt er. Før mente nordmenn at det er normalt å tvangssterilisere den samiske befolkningen, selv langt inn i de 20. århundre. Nå synes vi det er forkastelig. Hva har forandret seg? Dette er en reaksjon som kjennetegner \textit{moralsk realisme}
    \begin{quote}
        \textbf{\textit{Moralsk realisme:}} Det synet at det finnes objektive moralske fakta som er uavhengig av hva mennesker måtte synes eller føle om dem, og hva mennesker faktisk gjør.
    \end{quote}

    En motsatt reaksjon til oppdagelsen av vedvarende forskjeller mellom morale perspektiver er å ta disse som et tegn på at moralske verdier ikke egentlig eksisterer uavhengig av oss. Dette er reaksjon til \textit{moralsk antirealisme}

    \begin{quote}
        \textbf{\textit{Moralsk antirealisme:}} det synet at det ikke finnes noen objektive moralske sannheter eller verdier som ikke avhenger av hvordan vi mennesker tenker eller føler oss i forhold til dem eller andre skikker og praksiser.
    \end{quote}

    Vi kan forstå relativisme som en spesifikk type moralsk antirealsime. Relativisme er synet at det som er moralsk rett eller galt, alltid er relativt til en kultur eller historisk epoke. Hva som er rett og galt avhenger fra kulturer, normer og skikker. \newline

    En moralsk antirealist kan stille dette spørsmålet til en moralsk realist
    \begin{quote}
        Hvis ulike perioder og kulturer har ulike moralske oppfatninger, hvorfor tro at \textit{vi} har oppdaget den ene moralen som er sann for hele menneskeheten?
    \end{quote}
    Svaret vil kanskje da bli
    \begin{quote}
        Ville ikke også de andre kulturerene og historiske epokene si at vi tar feil? Eller at vi begge tar feil?
    \end{quote}
    De som hevder at moral kun er et spørsmål om synsing eller oppfatninger, er kulturrelativister. De mener at hva som er moralsk riktig avhenger av de moralske oppfatningene som dominerer i en kultur\newline

    Men antirealisme, dvs. å hevde at moralske verdier er sinnsavhengig, er ikke det samme som å akseptere kulturrelativistene. Kulturrelativisme er bare en version av antirealisme. Kan hende moralske verdier avhenger av aspekter ved menneskesinner som er de samme for alle menneskene. Tenk bare på Dawkins ide om egoistiske gener og hvordan de har formet menneskelig adgerd og menneskesinnet. Sharon Street foreslår at nettopp noen av våre mest grunnleggende moralske emosjoner er blitt formet av vår evolusjonshistorie, lenge før kulturer oppst. Hun argumenterer mot moralsk relativisme uten å argumentere for kulturrelativisme\newline


    Vi har nå \textit{moralsk epistemologi}. Her spør vi om hvordan vi får kunnskap om hva som er rett og galt, hvis det i det hele tatt er mulig. For eksempel forsøker Kant å utlede kategoriske imperativ fra bruk av fornuft alene. Moralsk kunnskap er ifølge Kant syntetisk a priori kunnskap. Denne kunnskap stammer ikke fra erfaring eller observasjon, men fra ren analyse av moralske begreper. Kants teori beskrives derfor ofte som et eksempel på \textit{rasjonalisme} om moralsk kunnskap. Rasjonalisme er ett eksempel på hvordan vi oppnår kunnskap om rett og galt.\newline

    Tenk på påstanden om at moralske oppfatninger og emosjoner varierer sterkt mellom ulike kulturer og epoker. Om vi spør en moralsk relativist hvordan vi kan vite at vi har rett hvis andre er uenige. Dette er et argument for at vi ikke kan vite om det er vi eller andre som har forstått de påstått objektive verdiene riktig. Synet at vi ikke vet eller kan vite hva som er moralsk rett eller galt, kalles moralsk skeptisisme. 
    \begin{quote}
        \textbf{\textit{Moralsk skeptisisme:}} Synet at vi ikke kan vite hva de virkelige moralske verdiene er. Det sier IKKE at det finnes objektive moralske verdier, bare at vi ikke kan erkjenne dem.
    \end{quote}

    I Sharon Street sin tekst, argumenterer hun for at antakelsen om at moralske verdier er helt uavhengige av oss, leder til moralsk skeptisisme. Hvis det fantes objektive moralske verdier kunne vi ikke vite hva de er. Dette er er grunn ifølge henne til å heller være moralsk antirealist.\newline


    Et tredje område av metaetikk handler om hva det er vi gjør når vi tenker eller sier at noen er moralske eller umoralske. Se for deg en person som mener moral er et spørsmål om oppfatning. De mener kanskje at det er umoralsk å rise barn, fordi det kan gi traumer. De uttrykker ikke et faktum, men en holdning. Dette kalles \textit{ekspressivisme}. Hvis ekspressivisme er sant, er det meningsløst å snakke om moralske objektive sannheter. Hvis dette er riktig, finnes det bokstavelig talt ikke oppfatninger om hva som er moralsk rett eller galt. (Wikforss kobler dette sammen mot å vite noe er sant og å tro på det)\newline

    Motparten til ekspressivisme er \textit{kognitivisme}. Det er det synet når vi sier noe er umoralsk, uttrykker vi en oppfatning om at noe er moralsk galt. Som nevnt handler Streets tekst om sammenheng mellom moralsk realisme (finnes objektive verdier?) og om hvordan moralsk kunnskap er mulig, hvis i det hele tatt. I resten av kapittelet antar vi at kognitivisme er riktig.
\end{flushleft}
\end{sectionbox}

\begin{sectionbox}[Genealogi]
\begin{flushleft}
    Vi undersøker nå hvordan personlig, kulturell og biologisk historie til moralske oppfatninger og følelser henger sammen med moralsk realisme og skeptisisme. Dette kalles for \textit{moralens genealogi}\newline

    Vi undersøker hvor våre moraler, normer og sannheter kommer fra, og vi viser at de ikke er evige, naturlige eller gitt. De har oppstått gjennom tilfeldigheter i historien, maktforhold og sosiale forhold

    \begin{quote}
        For eksempel: Mange mener at arbeid er moralsk godt. En genealogisk analyse vil spørre 
        \begin{enumerate}
            \item Når begynte vi å mene dette?
            \item Var det fordi arbeid er moralsk?
            \item Eller fordi makthavere trengte lydige arbeidere?
        \end{enumerate}
    \end{quote}
\end{flushleft}
\end{sectionbox}


\begin{sectionbox}[Genealogi, evolusjon og moral]
\begin{flushleft}
    Nietzche er en tysk filosof fra slutten av 1800-tallet. Han argumenterer for at når vi studerer moralens genealogi i europeisk historie, oppdager vi at vår samtidige idealer har sitt utspring i de middelmådige massenes forsøk på å kontrollere mektige individer. Moralske begreper er verktøy i en maktkamp mellom de svake massene og en mektig elite. Hans genealogi er av den undergravede sorten: Moralske begreper refererer ikke til objektive verdier\newline

    Han legger ikke frem mye evidens fordi det ikke fantes så mye av det på hans tid. Streets tekst gjør det, og bruker Darwins evolusjonsteori om naturlig utvalg. Spørsmålet Street tar for seg er hvorvidt det faktum at vi er resultatet av evolusjon via naturlig utvalg, undergraver ideen om at det finnes moralske sannheter, altså at den muliggjør en undergravende genealogi.\newline

    En enkel formulasjon av problemet er : \textit{Evolusjonsteori forklarer trekk ved organisme ut ifra hvorvidt trekke har bidratt til å øke organismenes fitness (antallet avkom det produserer). Men da må også våre moralske følelser og oppfatninger, i hvert fall delvis, forklares av det som økte vår evolusjonære fitness.}\newline


    \begin{quote}
        Darwin: \textit{"Om mennesker ble oppdratt under akkurat de samme forholdene som biene, kunne det ikke være noen tvil om at våre ugifte kvinner ville, som arbeiderne tenkt at deres hellige plikt er å drepe sine brødre, og mødre ville prøvd å drepe sine fruktbare døtre; og ingen kunne tenkt seg å gripe inn"}
    \end{quote}
    Ideen er at dersom vi hadde en annerledes evolusjonær historie, ville vi tenke veldig forskjellig om hva som er moralsk rett. Kant mener at vi jan utlede våre moralske plikter fra det kategoriske imperativ. Ville han utledet andre ting om han var en honningbie?
\end{flushleft}
\end{sectionbox}

\begin{sectionbox}[Studiespørsmål]
\begin{enumerate}
    \item Noen ganger kan det å lære årsaksopphavet til en holdning undergrave den. Andre ganger kan det å lære årsaksopphavet til en holdning styrke den. Street gir eksempler på begge typer effekter. Hva kunne være ekempler på en undergravende og en underbyggende genealogi?
    \item Hva skiller en sinnsuavhengig fra en sinnsavhengig forståelse av verdi?
    \item Hvorfor tenker Street at det evolusjonære opphavet til våre moralske følelser undergraver den sinnsuavhengige forståelsen av verdi? Hvordan kunne man svare på argumentene hennes?
    \item Hvorfor tenker Street at vi må utforske konkurrerende metaetiske synspunkter for å avgjøre hvorvidt evolusjonære forklaringer på våre mest grunnleggende verdier er en undergravende eller styrkende?
\end{enumerate}

\textbf{\textit{Dette bør du kunne}}
\begin{itemize}
    \item Forstå forskjellen mellom moralsk realisme og antirealisme
    \item Vite hva genealogi er, og hvorfor den kan være relevant for moralsk kunnskap
    \item Forklare hvorfor det faktum at vi er resultat av evolusjin, har noe å si for moralen
    \item Forstå Streets argument for moralsk antirealisme
\end{itemize}
\end{sectionbox}

\end{document}